\documentclass{beamer}

\mode<presentation>
{
  \usetheme[secheader]{Madrid}
  \setbeamercovered{transparent}
}

\usepackage[utf8]{inputenc}
\usepackage[spanish]{babel}
\usepackage{graphicx}
\usepackage{url}
\usepackage{times}
\usepackage[T1]{fontenc}
\usepackage{hyperref}
\hypersetup{pdfkeywords=open free source colombia Software Libre abierto simulación circuitos circuit simulation Freedom Day 2011 CIDCA AltaImpedancia alta impedancia electrónica gEDA gnucap qucs scilab}

\title[Simulación de circuitos con Software Libre]{Simulación de circuitos con Software Libre}

\author[JEGC]{Jorge~Ernesto~Guevara~Cuenca}

\institute[\url{www.altaimpedancia.org}]
{ \inst{1}
  Colibri - Comunidad de Usuarios de Software Libre en Colombia\\
  \url{http://www.slcolombia.org}
  \and
  \inst{2}
  Alta Impedancia\\
  \url{http://www.altaimpedancia.org}
  \and
  \inst{3}
  Hackbo - Hackerspace Bogotá\\
  \url{http://www.hackbo.co}
  \and
  \inst{4}
  Universidad Autónoma de Colombia\\
  \url{http://www.fuac.edu.co}
}

\date[SFD Bogotá 2011]
{Software Freedom Day Bogotá\\CIDCA\\17 de septiembre de 2011}

\subject{Simulación de circuitos con Sofware Libre, Software Freedom Day Bogotá, CIDCA, 17 de septiembre de 2011}

\logo{\includegraphics[height=1cm]{img/altaimpedancia}}

\AtBeginSection[]
{ \begin{frame}<beamer>{Agenda}
    \tableofcontents[currentsection,currentsubsection]
  \end{frame}}

\beamerdefaultoverlayspecification{<+->}

\begin{document}

\begin{frame}
  \titlepage
\end{frame}

\begin{frame}
  \frametitle{Agenda}
  \tableofcontents
\end{frame}

\section<presentation>*{Simulación de circuitos con Software Libre}

\begin{frame}{Objetivos}
  \begin{block}{Principal}
    Mostrar como simular circuitos electricos mediante ejemplos con los programas qucs, scilab y gEDA con gnucap.
  \end{block}
  \begin{block}{Especificos}
    \begin{itemize}
    \item Mostrar como usar los programas para hacer análisis transitorio y de frecuencia de algunos circuitos con dispositivos activos y pasivos.
    \item Hacer una comparación entre los programas mostrados durante la charla.
    \end{itemize}
  \end{block}
\end{frame}

\section{Terminología}

\begin{frame}
  \frametitle{Glosario}
  \begin{itemize}
  \item HDL acrónimo de Hardware Description Language (Lenguaje de descripción de hardware).
  \item PCB acrónimo de Printed Circuit Board (Tarjeta de circuito impreso).
  \item VLSI acrónimo de Very Large Scale Integration.
  \item FPGA acrónimo de Field Programable Gate Array.
  \item Netlist, lista de conexiones.
  \end{itemize}
\end{frame}

\begin{frame}{¿Qué es CAD?}
  \begin{itemize}
  \item CAD es el acrónimo de Computer Aided Design (Diseño asistido por computador).
  \item En principio se refiere a aplicaciones para hacer dibujos.
  \item Se usa en general para cualquier campo del conocimiento en el que se pueda diseñar por medio de un computador.
  \item ECAD es el acrónimo de Electronic Computer Aided Design (Diseño electrónico asistido por computador).
  \end{itemize}
\end{frame}

\begin{frame}{¿Qué es EDA?}
  \begin{itemize}
  \item EDA es el acrónimo de Electronic Design Automation (Automatización de diseño electrónico).
  \item Se refiere a todas las herramientas involucradas en el desarrollo de hardware.
  \item Es todo el software ECAD y hardware requeridos para el desarrollo de hardware.
  \item Hardware requerido
    \begin{itemize}
    \item HDK (Hardware Development Kits)
    \item HPP (Hardware Prototyping Plataforms).
    \end{itemize}
  \end{itemize}
\end{frame}

\begin{frame}{Herramientas EDA}
  \begin{itemize}
  \item Lenguajes de descripción de circuitos.
    \begin{itemize}
    \item netlist, ejemplo spice.
    \item HDL, ejemplo verilog.
    \end{itemize}
    \vspace{9pt}
  \item Captura esquemática y modelado gráfico.
    \begin{itemize}
    \item Diagramas de flujo.
    \item ASM.
    \item etc\dots{}
    \end{itemize}

    \vspace{9pt}
  \item Simuladores y visualizadores.
    \vspace{9pt}
  \item Fabricación de tarjetas de cirtuitos impresos (PCB).
  \item Fabricación de circuitos integrados (VLSI).
  \item Programación de dispositivos lógicos (FPGA, PAL, PLD).
  \end{itemize}
\end{frame}

\section{Presentación del software y primer ejemplo básico}

\subsection{qucs}

\begin{frame}
  
\end{frame}

\subsection{scilab}

\begin{frame}
  
\end{frame}

\subsection{gEDA}

\begin{frame}{gschem}{Capturador esquemático}
  \includegraphics[scale=0.5]{img/gschem.png}
\end{frame}

%Para configurar el fondo blanco en gschem se debe crear el archivo gschemrc en el directorio .gEDA dentro del home del usuario con el siguiente contenido:
% (load (build-path geda-rc-path "gschem-colormap-lightbg"))

\begin{frame}{gschem}{Instertar componentes}
  SPICE\\
  vdc-1\\
  gnucap-npn-1\\
\end{frame}

\begin{frame}[fragile]
  \framtitle{Ejecutar gschem}
  \begin{block}{}
\begin{verbatim}
gschem
\end{verbatim}
  \end{block}
\end{frame}

\begin{frame}
  gnetlist -s -g spice-sdb -o colpitts.ckt colpitts.sch
  gnucap -b colpitts.ckt
  gwave colpitts.dat
\end{frame}

\section{A donde ir por más...}

\begin{frame}{Otros programas}
  \begin{itemize}
  \item ktechlab...
  \end{itemize}
\end{frame}

\begin{frame}{Debian GNU/Linux}{¿Que tiene que ver el proyecto debian con esta charla?}
  \begin{itemize}
  \item Todo el software mencionado se encuentra empaquetado listo para usar
  \item También existe FEL
  \end{itemize}
\end{frame}

\section<presentation>*{Simulación de circuitos con Software Libre}

\begin{frame}{Conclusiones}
  \begin{block}{}
    \begin{itemize}
    \item El Diseño de hardware es un proceso que esta al alcance de todos, tiene grandes alcances y es importante buscar modelos que lleven el diseño a la producción en masa.
    \item Ya existen varios proyectos de gran envergadura que son Open Source Hardware.
    \item Hay un grave problema de interoperabilidad entre formatos y herramientas.
    \end{itemize}
  \end{block}
\end{frame}

\begin{frame}{Agradecimientos}
  \begin{itemize}
  \item Organización del Software Freedom Day.
  \item CIDCA
  \end{itemize}
\end{frame}

\appendix

\section<presentation>*{Referencias}

\begin{frame}
  \frametitle<presentation>{Bibliografía I}
  \begin{thebibliography}{99}
    \beamertemplatebookbibitems
  \bibitem[1]{PardoBoluda}Fernando Pardo Carpio y José A. Boluda Grau
    \newblock \emph{VHDL, Lenguaje para síntesis y modelado de circuitos, 2a. edición}
    \newblock Alfaomega, 2004
  \bibitem<1->[Bercelli 2006]{Bercelli}Ariel Bercelli
    \newblock \emph{Aprender la Libertad}
    \newblock 2006
    \newblock \url{http://www.aprenderlalibertad.org/libro-all}
  \end{thebibliography}
\end{frame}

\begin{frame}
  \frametitle<presentation>{Bibliografía II}
  \begin{thebibliography}{99}
    \beamertemplatearticlebibitems
  \bibitem<1->[Güichal 2005]{Guillermo}Guillermo Güichal
    \newblock \emph{Diseño Digital Utilizando Lógicas Programables}
    \newblock Junio 29, 2005
    \newblock \url{http://fpga.com.ar/notas/NotasCompletas.pdf}
  \bibitem<1->[4]{clasificacion}Ivan González, Juan González y Francisco Gómez-Arribas
    \newblock \emph{Hardware libre: clasificación y desarrollo de hardware reconfigurable en entornos GNU/Linux}
    \newblock 6 de Septiembre de 2003
    \newblock \url{http://www.iearobotics.com/personal/juan/publicaciones/art4/hardware-libre.pdf}
  \end{thebibliography}
\end{frame}

\begin{frame}[allowframebreaks]
  \frametitle<presentation>{Infografía}
  \begin{thebibliography}{99}
    \beamertemplatearticlebibitems
  \bibitem[5]{Ales}\emph{gEDA - GPL Electronic Design Automation}, \url{http://www.geda.seul.org/talks/deluge_ales.pdf}
  \bibitem[Oscilador colpitts]{colpitts}, \url{http://qucs.sourceforge.net/examples/colpitts_base.sch}
  \end{thebibliography}
\end{frame}

\section<presentation>*{Comunidades}

\begin{frame}{Como unirse}
  \begin{figure}
    \includegraphics{img/altaimpedancia-acerca-de}
  \end{figure}
  \begin{block}{}
    \begin{itemize}
    \item Linux en Caja - \url{http://linuxencaja.net}
    \item HackBo - \url{http://hackbo.co}
    \end{itemize}
  \end{block}
  \begin{block}{Contacto}
    \centering
    Jorge Ernesto Guevara Cuenca\\
    \href{mailto:ernesto@altaimpedancia.org}{ernesto@altaimpedancia.org}\\
  \end{block}
\end{frame}

\section<presentation>*{Sobre este documento}

\begin{frame}
  \begin{block}{Licencia}
    \begin{figure}
      \includegraphics[scale=0.9]{img/by-sa}
    \end{figure}
    \centering
    \small Creative Commons\\
    \small Atribución-Compartir Obras Derivadas Igual 2.5 Colombia\\
    \small \url{http://creativecommons.org/licenses/by-sa/2.5/co}
\end{block}
\begin{block}{}
  Creado con \LaTeX / Beamer
\end{block}
\end{frame}

\end{document}
